% Corrigé intégré automatiquement pour la banque Exercices.
% Source : B2_ProposerModele/B2_16_Hyperstatisme/69_TrainA350/69_TrainA350.tex
\begin{corrigebox}[Corrigé]
\CorrigeQuestion{1}
\begin{center}
\includegraphics[width=5cm]{69_01_cor}
\end{center}
\CorrigeQuestion{2}
La liaison de l'avion avec le sol est assimilabe à une liaison appui-plan de normale $\vect{z}$. Il y a donc 3 mobilités (1 rotation autour de $\vect{z}$, 1 translation selon $\vect{x}$ et 1 translation suivant $\vect{y}$.

En utilisant une méthode statique, on a $h=m-E_s+I_s$ avec : 
\begin{itemize}
\item  $m=3$;
\item $E_s = 1\times 6 = 6$ (on ne peut isoler que l'avion);
\item $I_s = 10 \times 1 = 10$ (8 liaisons ponctuelles avec 1 inconnue statique par liaison).
\end{itemize}
En conséquences, $h = 3-6+10 =7$. 

En utilisant une méthode cinématique, on a $h=m-I_c+E_c$ avec : 
\begin{itemize}
\item  $m=3$;
\item $E_c = \gamma \times 6 = (10-2+1) \times 6 = 54$ (on ne peut isoler que l'avion);
\item $I_c = 10 \times 5 = 50$ (8 liaisons ponctuelles avec 5 inconnues cinématiques par liaison).
\end{itemize}
En conséquences, $h = 3-50+54 =7$.
\CorrigeQuestion{3}
En utilisant une méthode statique, on a $h=m-E_s+I_s$ avec : 
\begin{itemize}
\item  $m=3$;
\item $E_s = 1\times 6 = 6$ (on ne peut isoler que l'avion);
\item $I_s = 3 \times 1 = 3$.
\end{itemize}
En conséquences, $h = 3-6+3 =0$. 

En utilisant une méthode cinématique, on a $h=m-I_c+E_c$ avec : 
\begin{itemize}
\item  $m=3$;
\item $E_c = \gamma \times 6 = (3-2+1) \times 6 = 12$ (on ne peut isoler que l'avion);
\item $I_c = 3 \times 5 = 15$ (3 liaisons ponctuelles avec 5 inconnues cinématiques par liaison);
\end{itemize}
En conséquences, $h = 3-15+12 =0$.
\end{corrigebox}
